% concatenated from main.tex
\documentclass[onefignum,onetabnum]{siamart220329}

\usepackage{amsmath,amssymb,amsfonts}
\usepackage{graphicx}
\usepackage{booktabs}
\usepackage{url}

% Compatibility patch for siamart220329/ntheorem with newer LaTeX kernels.
\makeatletter
\@ifundefined{cref@constructprefix}{}{%
  \def\refstepcounter@optarg[#1]#2{%
    \cref@old@refstepcounter{#2}%
    \cref@constructprefix{#2}{\cref@result}%
    \@ifundefined{cref@#1@alias}%
      {\def\@tempa{#1}}%
      {\def\@tempa{\csname cref@#1@alias\endcsname}}%
    \protected@edef\cref@currentlabel{%
      [\@tempa][\arabic{#2}][\cref@result]%
      \csname p@#2\endcsname\csname the#2\endcsname}%
  }%
}
\makeatother

\newcommand{\C}{\mathbb{C}}
\newcommand{\eps}{\varepsilon}
\newcommand{\dA}{\,dA}

\graphicspath{{figures/}{./}}

\begin{document}

\title{Cubature from rational approximation}

\author{Gentian Zavalani\thanks{Technische Universit{\"a}t Dresden, %
  Institute of Numerical Mathematics, %
  01062 Dresden, Germany
(\email{gentian.zavalani@tu-dresden.de}).}}

\headers{Cubature from rational approximation}{G.\ Zavalani}

\maketitle

\begin{abstract}
We present a numerical construction of cubature rules for area
integrals of analytic functions over planar domains with rectifiable
Jordan boundary. The starting point is the Cauchy--Green identity.
Given a weight $w$, we choose a $\bar\partial$-antiderivative $W$ and
reduce the area integral to a contour integral involving the boundary
values of $W$. These values are then approximated by a rational
function with free poles, computed by the AAA algorithm. The poles
inside the domain become cubature nodes, the corresponding residues
become weights, and the boundary residual controls the error through
an a posteriori estimate, rigorous once the continuous boundary
residual is bounded. The same rule admits a dual reading, as the exact
integral of a rational interpolant to the integrand, the area analogue
of the one-dimensional interpolatory viewpoint.
The numerical examples recover the disk mean-value rule and the
focal-segment rule of the ellipse to machine precision, reproduce the
exact finite quadrature identities of quadrature domains with both
separated and confluent nodes, and evaluate logarithmic and Cauchy
volume potentials from boundary data alone. The interior poles trace
analytic skeletons that we identify tentatively with the mother bodies
of potential theory, along with image points that appear without being
imposed; for the square the observed convergence is root-exponential.
\end{abstract}


\begin{keywords}
cubature, rational approximation, AAA algorithm, Cauchy--Green
formula,\linebreak
Schwarz function, quadrature domains, mother body, lightning
approximation, volume potentials
\end{keywords}


\begin{MSCcodes}
65D32, 41A20, 30E10, 31A05
\end{MSCcodes}

\section{Introduction}
\label{sec:intro}

In a recent paper, Horning and Trefethen showed how a great variety of
quadrature formulas in one dimension can be generated from rational
approximations of the Cauchy transform of a weight function
\cite{HorningTrefethen2026}. The poles of the rational function become
the quadrature nodes and the residues become the weights, and the whole
construction lives on arcs and contours in the complex plane. This
paper asks the corresponding two-dimensional question. Can quadrature
rules for \emph{area} integrals,
\begin{equation}
\label{eq:areaint}
I \;=\; \iint_\Omega f(z)\, w(z, \bar z) \dA ,
\end{equation}
over a domain $\Omega \subset \C$ with rectifiable Jordan boundary
$\gamma = \partial\Omega$, be generated in the same spirit? Throughout,
we assume that $f$ is analytic in a neighbourhood of $\bar\Omega$.

Remarkably, no new transform is required. A classical identity already
converts \eqref{eq:areaint} into a problem of approximation on the
boundary. Suppose $W(z,\bar z)$ satisfies
$\partial W/\partial\bar z=w$. Then the Cauchy--Green identity, which
underlies the Cauchy--Pompeiu formula \cite{Pompeiu}, gives, for $f$
analytic in a neighbourhood of $\bar\Omega$,
\begin{equation}
\label{eq:pompeiu}
\iint_\Omega f(z)\, w(z, \bar z) \dA
\;=\; \frac{1}{2i} \oint_\gamma f(z)\, W(z, \bar z)\, dz .
\end{equation}
The restriction of $W$ to $\gamma$ is known boundary data. In the
simplest case of a pure area integral we have $w \equiv 1$ and
$W = \bar z$, so the data is nothing but the complex conjugate of the
boundary parametrization. Now let $\rho$ be a rational function with
simple poles $z_k$ that approximates this data on $\gamma$. The poles
inside $\Omega$ contribute residues, while the poles outside $\Omega$,
together with any polynomial part, integrate to zero against analytic
$f$. Residue calculus thus turns \eqref{eq:pompeiu} into a
\emph{cubature rule},
\begin{equation}
\label{eq:rule}
I_n \;=\; \pi \sum_{z_k \in \Omega} c_k\, f(z_k),
\end{equation}
whose nodes are the interior poles and whose error is controlled
entirely by the quality of the fit on the boundary
(Theorem~\ref{thm:cert}). We compute the poles with the AAA algorithm
\cite{NST2018,NakatsukasaTrefethen2025,TrefethenWilber2025} and the
coefficients by a linear least-squares fit on the boundary samples. We
call \eqref{eq:rule} the \emph{AAA cubature rule}.

As in the one-dimensional case, the nodes and weights depend only on
$W$ and not on the integrand. Once $W|_\gamma$ has been approximated,
the same rule applies to every analytic $f$ for which
\eqref{eq:pompeiu} holds. The antiderivative is not unique: $W+h$
serves equally well for any analytic $h$, since
$\oint_\gamma f(z)h(z)\,dz=0$ leaves the exact contour integral
unchanged. At a finite rational degree, however, different choices of
$W$ lead to different numerical fits and hence to different computed
rules, and in each example below we simply use the natural choice,
stated explicitly. We make no claim to a general-purpose formula for
nonanalytic integrands, and the weights need not be positive or even
real.

It is worth emphasizing that the quantity entering the rigorous bound
lives entirely on the boundary: the residual in Theorem~\ref{thm:cert}
is $\|\rho-W\|_{\gamma,\infty}$, and nowhere do we assume anything about
$W$ in the interior of $\Omega$. The sampled validation residual used in
the computations is an a posteriori indicator of this quantity; it
becomes a certificate only once the continuous boundary supremum has
been bounded.

Where do the nodes end up? Their location is governed by the analytic
continuation of the boundary data into $\Omega$. When $w \equiv 1$ and
$\gamma$ is real-analytic, that continuation is the Schwarz function
$S(z)$ of the curve, the analytic function that matches $\bar z$ on
$\gamma$ \cite{Davis1974}. One therefore expects the near-best poles to
delineate the singularities of $S$ inside $\Omega$, tracing out what
potential theorists call a mother body, an analytic skeleton of the
domain \cite{Zidarov,Gustafsson1998,GustafssonPutinar}. Domains whose Schwarz function is
actually meromorphic in $\Omega$ are the \emph{quadrature domains} of
Aharonov and Shapiro
\cite{AharonovShapiro1976,Sakai1982,GustafssonShapiro2005}; these are
the domains that satisfy exact finite quadrature identities. Read in
this light, the construction is a direct numerical route from boundary
data to approximate quadrature identities. The inverse problem, that of
reconstructing a domain from its moments, has a substantial literature
\cite{GHMP2000,GPSS2009}, and one may regard what we do here as its
forward counterpart.

The work closest to ours is Trefethen's numerical computation of the
Schwarz function \cite{Trefethen2025Schwarz}, which applies AAA to
$\bar z$ on a curve and reads off branch cuts from the strings of poles
that appear. Horning and Trefethen \cite{HorningTrefethen2026}, as
already noted, obtain one-dimensional quadrature from rational
approximation of a Cauchy transform. The present paper joins these two
ideas through the Cauchy--Green identity: rational approximation of a
$\bar\partial$-antiderivative on the boundary produces a weighted
cubature rule whose error is controlled by the boundary residual.

Boundary reduction is of course not new. It underlies Gauss--Green
cubature \cite{SommarivaVianello2009}, which applies product Gauss
rules along the boundary, and connections between approximation and
analytic skeletons appear both in the zeros of Bergman polynomials
\cite{GPSS2009} and in explicit constructions of mother bodies for
special geometries \cite{Gustafsson1998,SavinaSterninShatalov2005}.
What rational approximation adds is a direct numerical construction that
attaches weights to the interior poles.

Throughout, we test the construction against exact identities,
independent reductions, adaptive area integration, and two-resolution
reference computations. Section~\ref{sec:classical} treats the
classical examples: it recovers the disk mean-value rule and the
focal-segment reduction of an ellipse, and it exhibits a candidate
five-armed skeleton for a starfish domain. Section~\ref{sec:qd} turns
to quadrature domains, where two examples separate the well-conditioned
case of distinct nodes from a confluent pair that represents a
derivative functional. Section~\ref{sec:square} takes up the square,
with its corner clustering and root-exponential convergence. Finally,
section~\ref{sec:potential} shows how the very same boundary
construction handles logarithmic and Cauchy volume potentials, for
targets both inside the domain and just outside it.


\section{The construction}
\label{sec:construction}

We parametrize $\gamma$ by $z(t)$, $t \in [0, 2\pi)$, positively
oriented, and write $\omega(t)=W(z(t),\overline{z(t)})$ for the boundary
data. We sample $\gamma$ at $M$ parameter values, taking care that the
grid resolves $\omega$; a near-boundary feature of scale $\delta$, in
particular, calls for spacing below $\delta$. Run to a prescribed
maximal degree with the sign-blended weight choice of
\cite{TrefethenWilber2025}, the option we use in all computations
reported here, AAA supplies a set of candidate poles. We discard the
nonfinite ones, and any pole within $10^{-8}d_M$ of a boundary sample,
where $d_M=\max_{j,k}|z_j-z_k|$ is the sampled boundary diameter. The
remaining poles are classified as interior or exterior with respect to
the sampled boundary, and we then solve the least-squares problem
\begin{equation}
\label{eq:ls}
\min_{c,\, d,\, p}\;
\sum_{j=1}^{M} \Bigl|
\sum_{z_k \in \Omega} \frac{c_k}{z_j - z_k}
+ \sum_{z_k \notin \Omega} \frac{d_k}{z_j - z_k}
+ \sum_{\ell=0}^{L} p_\ell\, z_j^\ell
- \omega_j \Bigr|^2 ,
\end{equation}
for a modest polynomial degree $L$. Although \eqref{eq:ls} is displayed
in an unscaled basis, the code works with $(z_j/s)^\ell$,
$s=\max_j|z_j|$, and rescales every column of the matrix to unit
2-norm. Writing $\sigma_1\geq\cdots$ for the singular values of this
scaled matrix, the rank-revealing solve keeps those with
\[
  \sigma_j>\tau, \qquad
  \tau=\max\{M,N\}\,\operatorname{eps}(\sigma_1),
\]
where $N$ is the number of columns and \texttt{eps} has its usual
MATLAB meaning, and the coefficients are returned in the original
partial-fraction basis. The exterior poles and the polynomial improve
the fit but integrate to zero against analytic integrands, so only the
interior coefficients survive into $I_n=\pi\sum c_kf(z_k)$. A pleasant
side benefit of fitting directly in this way is that we never have to
convert a clustered barycentric approximant into pole--residue form.

When a separate validation grid is available, we write $\eps_{\rm val}$
for its maximum residual. A convenient global measure of the
cancellation in the resulting rule is the quantity
\[
  \Lambda_n=\frac{\sum_k|\lambda_k|}
                   {\operatorname{area}(\Omega)},
  \qquad \lambda_k=\pi c_k.
\]
For a positive rule that reproduces the area we have $\Lambda_n=1$,
whereas large values warn of sensitivity to perturbations and of
cancellation. We stress that this is a diagnostic, not a stability
theorem.

In practice, when the exact integral is unknown, we raise the AAA
degree until $\eps_{\rm val}$ reaches the requested tolerance and stays
there as the validation grid is refined. For a given integrand,
$\tfrac12|\gamma|\eps_{\rm val}\max_\gamma|f|$ is then an a posteriori
error indicator, but not a rigorous bound, unless the continuous
residual itself has been bounded. We refine the fitting grid whenever
the fitting and validation residuals begin to separate, or a
near-boundary feature is left unresolved. And if $\Lambda_n$ grows
rapidly, or a confluent cluster appears, we take that as a signal that
the rule should be evaluated through the local moments of
section~\ref{sec:qd}; we claim no universal threshold for when this
happens.

\begin{theorem}\label{thm:cert}
Let $\Omega$ be a bounded domain whose boundary $\gamma$ is a
rectifiable Jordan curve, positively oriented, let $f$ be analytic in
a neighbourhood of $\bar\Omega$, and let $W$ be a
$\bar\partial$-antiderivative of $w$ for which the complex Green
identity \eqref{eq:pompeiu} holds. Let
\[
\rho(z)
= \sum_{z_k \in \Omega} \frac{c_k}{z-z_k}
+ \sum_{z_k \notin \Omega} \frac{d_k}{z-z_k}
+ p(z)
\]
be a rational approximant to $W$ on $\gamma$, where all poles are
simple, no pole lies on $\gamma$, and $p$ is a polynomial. Set
$\eps = \max_{z \in \gamma} |\rho(z) - W(z,\bar z)|$. Then the rule
$I_n = \pi \sum_{z_k \in \Omega} c_k f(z_k)$ satisfies
\[
|I - I_n| \;\le\; \tfrac{1}{2}\, |\gamma| \, \eps \,
\max_{z \in \gamma} |f(z)| ,
\]
where $|\gamma|$ is the length of $\gamma$.
\end{theorem}

\begin{proof}
On a rectifiable Jordan curve, Cauchy's theorem applies to functions
analytic in $\Omega$ and continuous on $\bar\Omega$ \cite{Walsh1933},
and the residue formula used next follows on subtracting from $f\rho$
its principal parts at the poles in $\Omega$.
Since $f$ is analytic in a neighbourhood of $\bar\Omega$ and the poles of $\rho$ avoid
$\gamma$, residue calculus gives
$\frac{1}{2i}\oint_\gamma f \rho \, dz = \pi \sum_{z_k \in \Omega}
c_k f(z_k) = I_n$: each interior pole contributes
$2\pi i\, c_k f(z_k)/2i$, and the exterior poles and the polynomial
part contribute zero. Subtracting this from \eqref{eq:pompeiu},
\[
I - I_n = \frac{1}{2i} \oint_\gamma f(z) \,
\bigl(W(z,\bar z) - \rho(z)\bigr)\, dz .
\]
Therefore
\[
|I-I_n|
\leq \frac12 \int_\gamma |f(z)|\, |W(z,\bar z)-\rho(z)|\, |dz|
\leq \tfrac12 |\gamma|\, \eps \max_{z\in\gamma}|f(z)|,
\]
which proves the estimate.
\end{proof}

Notice that Theorem~\ref{thm:cert} calls for the continuous supremum of
the residual on $\gamma$, which a finite sample does not deliver. A
fully rigorous certificate would require something more, such as an
adaptive or interval-based maximization along the boundary. What we can
say is that the independent validation-grid checks return residuals
consistent with those on the fitting grid. In fact the cubature error is
often much smaller than the indicator suggests, because the
contour-error integral enjoys a good deal of cancellation.


\paragraph{A dual interpolatory interpretation}
For pure area integrals there is a second way to look at the rule,
through the exterior area Cauchy transform
\[
   H(\xi) = \frac{1}{\pi}\iint_\Omega \frac{\dA(z)}{\xi-z},
   \qquad \xi \in \C \setminus \bar\Omega .
\]
This function is analytic outside $\bar\Omega$ and behaves like
$H(\xi)=\operatorname{area}(\Omega)/(\pi\xi)+O(\xi^{-2})$ at infinity.
The part of the fitted rational function that matters here is
\[
   r_n(z) = \sum_{z_k\in\Omega} \frac{c_k}{z-z_k}.
\]
Taking the Cauchy-kernel integrand $f_\xi(z)=1/(z-\xi)$ with $\xi$
outside $\bar\Omega$, the definition of $H$ and a direct evaluation of
the rule give
\begin{equation}
\label{eq:dual-cauchy}
   I(f_\xi) = -\pi H(\xi), \qquad
   I_n(f_\xi) = -\pi r_n(\xi), \qquad
   I(f_\xi)-I_n(f_\xi)=\pi\{r_n(\xi)-H(\xi)\}.
\end{equation}
In other words, the interior cubature rule is at the same time a
rational approximation to the exterior area Cauchy transform. For the
ellipse, where $H$ is known in closed form, one can check this identity
by hand.

The next statement makes this interpolatory reading precise. It is the
area analogue of the result in \cite[Section~11]{HorningTrefethen2026}.

\begin{proposition}[Conditional rational exactness]
\label{prop:dual}
Let $z_1,\ldots,z_n$ be the distinct interior cubature nodes, where
$n$ is their number, and set
\[
  r_n(z)=\sum_{k=1}^n\frac{c_k}{z-z_k},
  \qquad D(z)=r_n(z)-H(z).
\]
Suppose that $D$ has distinct zeros $s_1,\ldots,s_n$ in
$\C\setminus\bar\Omega$ and that the Cauchy matrix
\[
  C_{kj}=\frac{1}{z_k-s_j}
\]
is nonsingular. For any $f$ analytic in a neighbourhood of
$\bar\Omega$, let
\[
  q(z)=\sum_{j=1}^n\frac{\alpha_j}{z-s_j}
\]
be the unique rational function of this form satisfying
$q(z_k)=f(z_k)$. Then
\[
  I(q)=I_n(q)=I_n(f).
\]
\end{proposition}

\begin{proof}
Since the poles of $q$ lie outside $\bar\Omega$, Fubini's theorem and
the Cauchy formula give
\[
\frac{1}{2i}\oint_\gamma q(\xi)H(\xi)\,d\xi
 =\frac{1}{2\pi i}\iint_\Omega
   \oint_\gamma \frac{q(\xi)}{\xi-z}\,d\xi\,dA(z)
 =\iint_\Omega q(z)\,dA(z).
\]
Residue calculus also gives
$I_n(q)=(2i)^{-1}\oint_\gamma q(z)r_n(z)\,dz$. Hence
\[
I(q)-I_n(q)
  =-\frac{1}{2i}\oint_\gamma q(z)D(z)\,dz.
\]
The exterior poles of $q$ are cancelled by the zeros of $D$, so $qD$
is analytic in the exterior of $\gamma$. Moreover,
\[
  q(z)=O(z^{-1}),\qquad D(z)=O(z^{-1}),
\]
and therefore $q(z)D(z)=O(z^{-2})$. Deforming $\gamma$ to a large
circle in the exterior shows that the large-contour contribution
vanishes. Thus
$I(q)=I_n(q)$, while the interpolation conditions give
$I_n(q)=I_n(f)$.
\end{proof}

The reading applies only when the effective residual happens to have
enough suitable exterior zeros and the associated Cauchy matrix is
nonsingular. Neither condition is guaranteed by the least-squares
construction, and in practice we identify the zeros a posteriori. The
exterior-pole and polynomial terms in \eqref{eq:ls} are auxiliary ---
they leave $I_n$ unchanged for analytic integrands --- which is why the
proposition works with $r_n$ alone; were one to use the full boundary
fit in an exterior deformation, the interpolant would have to cancel
these auxiliary poles as well. Expanding \eqref{eq:dual-cauchy} at
infinity, finally, gives
\[
   D(z)=\frac{\sum_k c_k-\operatorname{area}(\Omega)/\pi}{z}
        +O(z^{-2}),
\]
so that exact reproduction of the area improves $D$ to $O(z^{-2})$ and
allows broader normalizations at infinity. This is not needed for the
strictly proper $q$ of Proposition~\ref{prop:dual}, for which
$qD=O(z^{-2})$ already holds. The whole conditional interpretation
carries over to other weights, with $H$ replaced by the exterior Cauchy
transform of $w\dA$.

In our experience the results are insensitive to modest changes in the
polynomial degree $L$. The full auxiliary basis can nonetheless be
numerically redundant: a remote exterior pole, for instance, is
effectively a pole at infinity, and its column may duplicate part of
the polynomial span. The rank-revealing solve removes such redundancy
while leaving the interior-pole rule untouched, and so the conditioning
of the auxiliary fit need not say anything about the conditioning of the
cubature rule itself.

There are two practical points worth stating explicitly. First, the identity
\eqref{eq:pompeiu} needs $W$ to be single-valued on $\bar\Omega$ and
smooth apart from integrable singularities; the antiderivatives of
section~\ref{sec:potential} are chosen to be single-valued and bounded
at the target, so that no correction terms are needed whether the
target lies inside or outside $\Omega$. Second, with $\gamma$
positively oriented the factor $\pi$ in \eqref{eq:rule} is just
$2\pi i/(2i)$, and the rule reproduces $\mathrm{area}(\Omega)$ when
$f \equiv 1$ --- a convenient check to run at execution time. One must
still watch for Froissart-type pole--zero pairs. A spurious interior
pole with a small least-squares coefficient may do no harm, though this
is not guaranteed in an ill-conditioned basis; a confluent pair with
large opposite weights, on the other hand, can be entirely genuine and
is best handled through the local moments of section~\ref{sec:qd}.

Beyond the ability to evaluate $W$ on the boundary, nothing else in the
construction depends on the weight at all; section~\ref{sec:potential}
makes use of two different $\bar\partial$-antiderivatives for potential
kernels.

\section{Classical examples}
\label{sec:classical}

Unless we say otherwise, relative error means $|I_n-I|/|I|$. When a
maximum is taken over several integrands, we switch to the scaled error
$|I_n-I|/\max(1,|I|)$, so that a zero or very small reference value
cannot distort the comparison. Table~\ref{tab:summary} collects the
smooth, quadrature-domain, and polygonal examples of this and the next
two sections; the volume potentials of
section~\ref{sec:potential} have their own tables.

\begin{table}[t]
\centering
\caption{The smooth, quadrature-domain, and polygonal examples: the
number $n$ of nodes carrying the rule and the best relative error
(worst case over the test integrands, where several are used), at the
degrees used in the corresponding experiments.}
\label{tab:summary}
\begin{tabular}{lcc}
\toprule
example & $n$ & rel.\ err.\\
\midrule
disk, mean-value rule                     & 1   & $3.6\times10^{-15}$\\
ellipse, focal Gauss--Gegenbauer          & 38  & $8.0\times10^{-16}$\\
starfish, adaptive two-dimensional reference & 55  & $4.1\times10^{-16}$\\
quadrature domain, confluent pair         & 2   & $8.6\times10^{-11}$\\
Neumann's oval, separated nodes           & 2   & $8.8\times10^{-15}$\\
square, closed form / tensor Gauss        & 159 & $1.6\times10^{-11}$\\
\bottomrule
\end{tabular}
\end{table}

\paragraph{Disk} On a circle of radius $R$ centered at $z_c$ the
boundary data is already rational, since
$\bar z = \bar z_c + R^2/(z - z_c)$ on $\gamma$. The only non-negligible
interior contribution collapses onto the center, with weight $\pi R^2$,
while whatever additional poles the algorithm produces carry negligible
weight. Rational approximation of the boundary data thus reproduces the
classical mean value property, exactly as one would hope.

\begin{figure}[t]
\centering
\includegraphics[width=\textwidth]{cub2_ellipse.pdf}
\caption{Ellipse with semiaxes $(1, 0.6)$. Left: the AAA cubature
nodes (dots) cluster near the focal segment; open circles are the
Gauss--Gegenbauer points of the exact reduced rule; crosses mark the
foci. Right: relative error of the AAA cubature rule and of
Gauss--Gegenbauer on the focal segment, for $f = e^z$.}
\label{fig:ellipse}
\end{figure}

\paragraph{Ellipse} Take the ellipse with semiaxes $(a, b) = (1, 0.6)$
and $c^2 = a^2 - b^2$. Here the Schwarz function has a branch cut along
the focal segment $[-c, c]$, and shrinking the contour of
\eqref{eq:pompeiu} down onto that cut gives the exact reduction
\begin{equation}
\label{eq:focal}
\iint_E f \dA \;=\; \frac{2ab}{c^2} \int_{-c}^{c} f(x)\,
\sqrt{c^2 - x^2}\, dx ,
\end{equation}
a Gegenbauer (Chebyshev second-kind) weight on the focal segment. We
verified \eqref{eq:focal} independently. As Figure~\ref{fig:ellipse}
shows, the AAA nodes cluster near the focal segment even though the
algorithm is never told where that segment is, and the figure compares
their convergence with Gauss--Gegenbauer quadrature.

This comparison deserves a word of interpretation. Gauss--Gegenbauer
converges faster, as one would expect; its rate is set by analyticity
relative to the focal segment and is superexponential for $e^z$. The
AAA nodes and weights, by contrast, are built with no knowledge of $f$,
and their residual-controlled worst-case behavior reflects rational
approximation of the boundary data. The error for any one integrand
depends further on its size on the boundary and on its analytic
continuation, and can be far smaller than the worst case thanks to
cancellation in the contour-error integral. It is worth recalling that
in the one-dimensional near-best setting of \cite{HorningTrefethen2026},
Gauss-type rules carry twice the exponent of the corresponding rational
construction --- the Gauss-quadrature factor of two. So when the exact
reduction \eqref{eq:focal} happens to be known, the classical rule
holds the advantage.

But what if the geometric information is imperfect? Figure~\ref{fig:perturbed}
shows the answer. Gauss--Gegenbauer on the focal segment of the
unperturbed ellipse stalls at the level $O(p)$ of the geometry error,
whereas the AAA rule, working from the true boundary data, drives on
down to machine precision. The classical rule is only as accurate as
its geometric model; the AAA rule adapts to the perturbed boundary as
it actually is.

\begin{figure}[t]
\centering
\includegraphics[width=\textwidth]{cub6_perturbed.pdf}
\caption{The ellipse perturbed by $p\, e^{3it}$, with $p = 10^{-2}$
(left) and $p = 10^{-4}$ (right); function $f = e^z$.
Gauss--Gegenbauer on the focal segment of the underlying ellipse
floors at the geometry error $O(p)$ (dotted line at height $p$); the
AAA cubature rule, built from the true boundary data, converges to
machine precision.}
\label{fig:perturbed}
\end{figure}

\begin{figure}[t]
\centering
\includegraphics[width=\textwidth]{cub3_domains.pdf}
\caption{Starfish domain $r(t) = 1 + 0.3\cos 5t$. Left: the cubature
nodes trace a five-armed mother-body-like analytic skeleton of the
domain. Right: relative error of the AAA cubature rule for an entire
integrand and for an integrand with a pole at $z = 1.6$.}
\label{fig:starfish}
\end{figure}

\paragraph{Starfish} For the domain $r(t) = 1 + 0.3 \cos 5t$ the
Schwarz function is no longer elementary, and its interior
singularities weave a more intricate skeleton. The outcome is shown in
Figure~\ref{fig:starfish}: the nodes trace a five-armed star inside the
domain, a candidate mother-body-like skeleton of $\Omega$
\cite{Zidarov,Gustafsson1998}, and the rule converges geometrically all
the way to machine precision. As an independent check, the boundary
reference agrees with two-dimensional adaptive integration over
$\Omega$, carried out separately on the real and imaginary parts.

\section{Quadrature domains}
\label{sec:qd}

Recall that a domain $\Omega$ is a quadrature domain if there are
finitely many points $a_j \in \Omega$ and coefficients $c_{jk}$, only
finitely many of them nonzero, for which
\[
\iint_\Omega f \dA \;=\; \sum_j \sum_{k \ge 0} c_{jk} f^{(k)}(a_j)
\qquad \text{for all integrable analytic } f,
\]
or, equivalently, if its Schwarz function is meromorphic in $\Omega$
\cite{AharonovShapiro1976,Sakai1982,GustafssonShapiro2005}. The
simplest nontrivial example is $\Omega = \phi(\mathbb{D})$ with
$\phi(\zeta) = \zeta + a\zeta^2$, $0 < a < 1/2$, whose Schwarz
function has a single double pole at $0$ and therefore obeys the exact
identity
\begin{equation}
\label{eq:asidentity}
\iint_\Omega f \dA \;=\; \pi \bigl[ (1 + 2a^2)\, f(0) + a\, f'(0) \bigr].
\end{equation}
We confirmed \eqref{eq:asidentity} against an independent boundary
reference.

\begin{figure}[t]
\centering
\includegraphics[width=\textwidth]{cub3b_qd.pdf}
\caption{The quadrature domain $\phi(\mathbb{D})$,
$\phi(\zeta) = \zeta + 0.3\zeta^2$. Left: the discovered rule consists
of two nodes $1.4 \times 10^{-5}$ apart with large opposite weights, a
confluent pair representing the node of order two at the origin in
\eqref{eq:asidentity}. Right: relative error against AAA degree; the
dominant two-node confluent structure first appears at degree $2$ and
persists across the displayed range; the right panel shows the
corresponding integration error.}
\label{fig:qd}
\end{figure}

How does a rule with simple poles reproduce a derivative functional?
The AAA rule recovers \eqref{eq:asidentity} through what we call a
\emph{confluent pair}: since the fit uses only simple poles, the double
pole of the Schwarz function is mimicked by two nearby nodes carrying
large, nearly opposite weights. It is the zeroth and first local
moments that recover the two coefficients of the exact identity; the
individual nodes and weights are not the stable quantities.

The representation is accurate, if less stable than a rule with
well-separated nodes, and the large opposite weights are nothing but a
finite-difference stand-in for the derivative functional in
\eqref{eq:asidentity}. For a cluster $C$ with center $z_C$, the
quantities that are actually stable are the local moments
\[
   \mu_j = \sum_{k\in C} \lambda_k (z_k-z_C)^j,
\]
not the individual weights. Expanding $f$ about $z_C$,
\[
   \sum_{k\in C} \lambda_k f(z_k)
   = \sum_{j=0}^{r-1} \frac{\mu_j}{j!} f^{(j)}(z_C)
     + O\!\left(\max_{z\in C}|f^{(r)}(z)|
       \sum_{k\in C}|\lambda_k|\,|z_k-z_C|^r\right),
\]
one sees that a nearly confluent simple-pole rule ought to be
interpreted, and where possible evaluated, in this derivative-node
basis. A practical, if heuristic, compression rests on the two
scale-invariant diagnostics
\[
\eta_{\rm sep}(C)
 =\frac{\max_{j,k\in C}|z_j-z_k|}
        {\operatorname{diam}(\Omega)},
\qquad
\eta_{\rm cancel}(C)
 =\frac{|\sum_{k\in C}\lambda_k|}
        {\sum_{k\in C}|\lambda_k|}.
\]
We flag a cluster when both of these are small and its share of
$\Lambda_n$ is large. A production code would need to fix and test
thresholds for the diagnostics; here they serve only to guide
interpretation. Once a cluster has been spotted we set
\[
   z_C=\frac{\sum_{k\in C}|\lambda_k|z_k}
             {\sum_{k\in C}|\lambda_k|},
\]
and evaluate it through the moments $\mu_j$ rather than by summing the
raw weights directly. For a general smooth domain the same computation
returns an approximate quadrature identity together with a validation
residual, a forward complement to the inverse shape-from-moments
problem \cite{GHMP2000}.

\begin{figure}[t]
\centering
\includegraphics[width=\textwidth]{cub9_neumann.pdf}
\caption{Neumann's oval, $\phi(\zeta) = \zeta/(1 - 0.25\zeta^2)$.
Left: the two computed nodes (dots) and the exact nodes $\pm 8/15$ of
\eqref{eq:neumann} (crosses). Right: maximal scaled error over five
integrands against the AAA degree; the number of poles placed
inside $\Omega$ is exactly two at every degree.}
\label{fig:neumann}
\end{figure}

A quadrature domain with \emph{separated} simple nodes behaves better
still. Neumann's oval is $\Omega = \phi(\mathbb{D})$ with
$\phi(\zeta) = \zeta/(1 - q^2 \zeta^2)$ and $0 < q < 1$
\cite{GustafssonShapiro2005}. This map is univalent on $\mathbb{D}$,
its Schwarz function has two simple poles, and residue calculus on the
unit circle yields the exact identity
\begin{equation}
\label{eq:neumann}
\iint_\Omega f \dA \;=\; \pi c \,\bigl[ f(x_*) + f(-x_*) \bigr],
\qquad
x_* = \frac{q}{1 - q^4}, \quad
c = \frac{1 + q^4}{2\,(1 - q^4)^2} ,
\end{equation}
with both nodes $\pm x_*$ lying interior to $\Omega$.
Figure~\ref{fig:neumann} treats the case $q=1/2$, for which $x_*=8/15$
and $c=136/225$. The two nodes come out correctly at every degree
shown, and raising the degree merely sharpens the exterior
approximation, driving the integration error to machine precision. This
separated rule is both more accurate and better conditioned than the
confluent one above.

\section{A domain with corners}
\label{sec:square}

The error theorem holds for any rectifiable Jordan curve, and the
experiments in this paper all use smooth or piecewise-analytic
boundaries, with parametrizations and sampling grids fine enough to
resolve the boundary data.

At a corner, though, the character of the approximation problem
changes. Along a straight edge through $a$ in the direction
$e^{i\theta}$ we have $\bar z=\bar a+e^{-2i\theta}(z-a)$, but two
adjacent edges continue to different functions. The singularities
thrown up at the corners bring exponential clustering of the poles and
root-exponential convergence $O(e^{-C\sqrt n})$
\cite{GopalTrefethen2019} --- the familiar signature of lightning
approximation. For the square the observed nodes line up along the two
diagonals, just as a mother-body configuration would suggest
\cite{Gustafsson1998}.

Figure~\ref{fig:square} shows the computation for the square
$[-1,1]^2$, with the boundary sampling graded toward the corners. The
nodes cluster at the corners and hug the diagonals, while the exterior
poles run off along the corresponding corner rays. The closed-form
reference for $e^z$ and an independent tensor Gauss--Legendre reference
agree to machine precision. As the figure shows, the cubature error
falls root-exponentially, exactly as lightning approximation predicts.

\begin{figure}[t]
\centering
\includegraphics[width=\textwidth]{cub8_square.pdf}
\caption{The square $[-1,1]^2$. Left: cubature nodes (dots) align with
the two diagonal candidate mother-body segments, with exponential
clustering toward the corners; the gray dots are exterior poles, which
also line up on the corner rays of the exterior skeleton. Right: relative
error against $\sqrt{n}$ for two integrands; the dashed line is
the least-squares fit $\propto e^{-1.90\sqrt{n}}$, the
root-exponential regime of lightning approximation.}
\label{fig:square}
\end{figure}

\section{Nearly singular volume potentials}
\label{sec:potential}

As a last application, consider the two-dimensional Newtonian potential
of an analytic integrand $f$ over $\Omega$,
\begin{equation}
\label{eq:newton}
\Phi(z_0) \;=\; \iint_\Omega f(z)\, \log|z - z_0| \dA ,
\end{equation}
which is singular when $z_0\in\Omega$ and nearly singular when $z_0$
sits close to $\gamma$. For general densities one commonly resorts to
volume meshes with local corrections, or to density interpolation
\cite{AndersonEtAl2024,ShenSerkh2024}; for analytic $f$, however, the
present construction again needs the boundary data alone. The
$\bar\partial$-antiderivative
\begin{equation}
\label{eq:Wlog}
W(z, \bar z) \;=\; (\bar z - \bar z_0)\bigl( \log|z - z_0| -
\tfrac{1}{2} \bigr)
\end{equation}
is single-valued and continuous for $z_0$ inside or outside $\Omega$,
provided only that $z_0\notin\gamma$, and it treats interior and
exterior targets on the same footing. Near $\gamma$ the sampling must
of course resolve a boundary feature on the scale of the target
distance. For $z_0\in\Omega$ we apply \eqref{eq:pompeiu} after excising
a disk of radius $\epsilon$ about $z_0$; since
$|W|=O(\epsilon\log\epsilon)$ on the small circle, its contribution is
$O(\epsilon^2\log\epsilon)$ and disappears as $\epsilon\to0$, the
logarithmic kernel being locally integrable. The same argument takes
care of the bounded Cauchy antiderivative introduced below.

\begin{figure}[t]
\centering
\includegraphics[width=\textwidth]{cub4_potential.pdf}
\caption{Newtonian potential over the starfish body. Left: nodes for
an interior target near the boundary (red star: the nodes cluster at
the target) and for an exterior target at distance $0.05$ (black star:
the nodes cluster at a candidate reflected point inside the body); both
node sets also trace the mother-body-like analytic skeleton. Right: relative error against
the number of nodes for four targets; the exterior case at distance
$0.01$ is limited here by the $M = 800$ boundary sampling (see
Table~\ref{tab:resM}).}
\label{fig:potential}
\end{figure}

\begin{table}[t]
\centering
\caption{Volume potential \eqref{eq:newton} over the starfish body:
best relative errors for four targets, with $M = 800$ boundary
samples.}
\label{tab:res}
\begin{tabular}{lcc}
\toprule
target & rel.\ err. & $n$\\
\midrule
interior, central        & $2.2\times 10^{-16}$ & 67\\
interior, near boundary  & $6.5\times 10^{-16}$ & 64\\
exterior, $\delta=0.05$  & $8.5\times 10^{-16}$ & 61\\
exterior, $\delta=0.01$  & $6.2\times 10^{-10}$ & 68\\
\bottomrule
\end{tabular}
\end{table}

\begin{table}[t]
\centering
\caption{The exterior target at distance $\delta = 0.01$ as the
boundary sampling is refined: the accuracy floor belongs to the
sampling, not to the method.}
\label{tab:resM}
\begin{tabular}{ccc}
\toprule
$M$ & rel.\ err. & $n$\\
\midrule
800  & $6.2\times 10^{-10}$ & 68\\
1600 & $3.0\times 10^{-14}$ & 68\\
3200 & $2.2\times 10^{-16}$ & 63\\
\bottomrule
\end{tabular}
\end{table}

Figure~\ref{fig:potential} and Tables~\ref{tab:res} and~\ref{tab:resM}
collect the results. Every well-resolved target reaches machine
precision, and for the closest exterior target, refining the boundary
grid removes what at first looks like an accuracy floor.

The geometry of the nodes is itself informative. For an interior target
the boundary data \eqref{eq:Wlog} continues inward with a branch point
at $z_0$, and the nodes duly cluster at the target. For an exterior
target they cluster instead near a point consistent with a solution of
$S(z)=\bar z_0$ --- the relevant reflected branch in this experiment ---
which for a circle is exactly the classical inverse point of Kelvin. So
the left panel of Figure~\ref{fig:potential} offers numerical evidence
of an image-point structure that was in no way imposed in advance.

The construction is kernel-independent in a precise sense: only the
boundary antiderivative changes. Replacing \eqref{eq:Wlog} by
$W = (\bar z - \bar z_0)/(z - z_0)$, which is bounded at $z_0$, produces
the area Cauchy integral $\iint_\Omega f(z)/(z - z_0) \dA$. Nothing else
in the construction is altered, and once more the numerical check comes
in near machine precision.

This same area Cauchy integral delivers a derivative of the potential.
For $z_0$ outside $\bar\Omega$, ordinary differentiation under the
integral sign applies. For $z_0\in\Omega$, one may split off a disk
centered at $z_0$ and pass to the limit, since $(z-z_0)^{-1}$ is locally
integrable. As
\[
  \frac{\partial}{\partial z_0}\log|z-z_0|
  =-\frac{1}{2(z-z_0)},
\]
we obtain
\[
\iint_\Omega \frac{f(z)}{z - z_0} \dA \;=\;
-2\, \frac{\partial \Phi}{\partial z_0} ,
\]
where $\partial/\partial z_0 = \tfrac12(\partial/
\partial x_0-i\,\partial/\partial y_0)$. For a real-valued potential
this reproduces the usual gradient components; for complex $f$ it is to
be read componentwise.

\section{Discussion}
\label{sec:discussion}

We have followed a single chain of reasoning from the Cauchy--Green
identity to boundary rational approximation, to interior poles and
residues, and finally to cubature. The error theorem covers analytic
integrands and a fixed admissible $\bar\partial$-antiderivative; it says
nothing about general continuous integrands. The weights may be complex
or signed, the sampled residuals stay mere indicators until the
continuous residual has been bounded, and confluent clusters have to be
evaluated through their moments.

The main open questions all concern the geometry of the nodes. Our
starfish and potential experiments give numerical evidence of a
candidate mother-body-like skeleton and of reflected image points, but
they do not prove that the displayed sets are the exact objects of
potential theory. On the square the measured rate is consistent with
the root-exponential behavior of lightning approximation
\cite{GopalTrefethen2019}, yet a proof in this setting is another
matter. For the derivative identities, a backward-error analysis of the
moment compression is still missing. And although in the dual
experiment a rank-revealing solve strips away a redundant auxiliary
direction without disturbing the interior-pole rule, this is encouraging
evidence rather than a uniform stability theorem. Each of these is an
invitation for further work.

Several extensions suggest themselves. The boundary identity carries
over directly to domains with several boundary components, though we
have not tested the resulting pole classification here. General
polygons and curved corners will demand the same care with graded
sampling as the square did. Other kernels will call for their own
$\bar\partial$-antiderivatives, and families of targets for a
set-valued AAA approximation with shared poles. A genuinely
three-dimensional analog would need some replacement for the planar
residue calculus, and that lies well beyond the present paper.

All results were generated using MATLAB based on Chebfun's \texttt{aaa}. The scripts are publicly available at \url{https://github.com/zavala92/cubature_aaa}.

\section*{Acknowledgments}

We are grateful to Nick Trefethen for helpful comments on an earlier
draft, particularly concerning the analyticity assumptions and the role of
boundary approximation, and to Andrew Horning for discussions of the
dual interpretation and numerical stability.


\begin{thebibliography}{99}
\itemsep=1pt

\bibitem{AharonovShapiro1976}
D.~Aharonov and H.~S. Shapiro,
\emph{Domains on which analytic functions satisfy quadrature
identities}, J. Analyse Math., 30 (1976), pp.~39--73.

\bibitem{AndersonEtAl2024}
T.~G. Anderson, M.~Bonnet, L.~M. Faria, and C.~P\'erez-Arancibia,
\emph{Fast, high-order numerical evaluation of volume potentials via
polynomial density interpolation}, J. Comput. Phys., 511 (2024),
113161.

\bibitem{Davis1974}
P.~J. Davis,
\emph{The Schwarz Function and its Applications}, Carus Mathematical
Monographs 17, Mathematical Association of America, 1974.

\bibitem{GopalTrefethen2019}
A.~Gopal and L.~N. Trefethen,
\emph{Solving Laplace problems with corner singularities via rational
functions}, SIAM J. Numer. Anal., 57 (2019), pp.~2074--2094.

\bibitem{Gustafsson1998}
B.~Gustafsson,
\emph{On mother bodies of convex polyhedra}, SIAM J. Math. Anal., 29
(1998), pp.~1106--1117.

\bibitem{GHMP2000}
B.~Gustafsson, C.~He, P.~Milanfar, and M.~Putinar,
\emph{Reconstructing planar domains from their moments},
Inverse Problems, 16 (2000), pp.~1053--1070.

\bibitem{GustafssonPutinar}
B.~Gustafsson and M.~Putinar,
\emph{Hyponormal Quantization of Planar Domains}, Lecture Notes in
Mathematics 2199, Springer, 2017.

\bibitem{GPSS2009}
B.~Gustafsson, M.~Putinar, E.~B. Saff, and N.~Stylianopoulos,
\emph{Bergman polynomials on an archipelago: estimates, zeros and
shape reconstruction}, Adv. Math., 222 (2009), pp.~1405--1460.

\bibitem{GustafssonShapiro2005}
B.~Gustafsson and H.~S. Shapiro,
\emph{What is a quadrature domain?}, in Quadrature Domains and Their
Applications, Oper. Theory Adv. Appl. 156, Birkh\"auser, 2005,
pp.~1--25.

\bibitem{HorningTrefethen2026}
A.~Horning and L.~N. Trefethen,
\emph{Quadrature formulas from rational approximations},
IMA J. Numer. Anal. (2026), doi:10.1093/imanum/draf138.

\bibitem{NST2018}
Y.~Nakatsukasa, O.~S\`ete, and L.~N. Trefethen,
\emph{The AAA algorithm for rational approximation}, SIAM J. Sci.
Comput., 40 (2018), pp.~A1494--A1522.

\bibitem{NakatsukasaTrefethen2025}
Y.~Nakatsukasa and L.~N. Trefethen,
\emph{Applications of AAA rational approximation},
Acta Numerica, 35 (2026), pp.~459--601.

\bibitem{Pompeiu}
D.~Pompeiu,
\emph{Sur une classe de fonctions d'une variable complexe},
Rend. Circ. Mat. Palermo, 33 (1912), pp.~108--113.

\bibitem{Sakai1982}
M.~Sakai,
\emph{Quadrature Domains}, Lecture Notes in Mathematics 934,
Springer, 1982.

\bibitem{SavinaSterninShatalov2005}
T.~V. Savina, B.~Yu. Sternin, and V.~E. Shatalov,
\emph{On a minimal element for a family of bodies producing the same
external gravitational field}, Appl. Anal., 84 (2005), pp.~649--668.

\bibitem{ShenSerkh2024}
Z.~Shen and K.~Serkh,
\emph{Rapid evaluation of Newtonian potentials on planar domains},
SIAM J. Sci. Comput., 46 (2024), pp.~A609--A628.

\bibitem{SommarivaVianello2009}
A.~Sommariva and M.~Vianello,
\emph{Gauss--Green cubature and moment computation over arbitrary
geometries}, J. Comput. Appl. Math., 231 (2009), pp.~886--896.

\bibitem{Trefethen2025Schwarz}
L.~N. Trefethen,
\emph{Numerical computation of the Schwarz function},
Constructive Approximation (2025), doi:10.1007/s00365-025-09719-2.

\bibitem{TrefethenWilber2025}
L.~N. Trefethen and H.~D. Wilber,
\emph{Computation of Zolotarev rational functions}, SIAM J. Sci.
Comput., 47 (2025), pp.~A2205--A2220.

\bibitem{Walsh1933}
J.~L. Walsh,
\emph{The Cauchy--Goursat theorem for rectifiable Jordan curves},
Proc. Natl. Acad. Sci. USA, 19 (1933), pp.~540--541.

\bibitem{Zidarov}
D.~Zidarov,
\emph{Inverse Gravimetric Problem in Geoprospecting and Geodesy},
Elsevier, Amsterdam, 1990.

\end{thebibliography}

\end{document}
